%% kapa: exam=1 round=37 year=2026 period=1 form=A subject=economics subject_ko=경제학원론 seq=40 number=80
%% source: Q-Net 감정평가사 기출문제 (cst00309) artlSeq=5258924 file=raw/1008/5258924/2026년 제37회 감정평가사 1차 1교시.pdf
%% conversion: kordoc 4.24.51 -> tools/gen1.py (auto); review=proofread-sonnet; flags=none
\begin{question}{80}
1년 간 물가상승률이 한국은 3 \%, 미국은 2 \%가 될 것으로 예상되며, 한국의 1년 만기 명목이자율은 5 \%이다. 구매력평가설(purchasing power parity)과 이자율평가설(interest rate parity)이 성립하는 경우, 이에 관한 설명으로 옳은 것을 모두 고른 것은? (단, 원-달러 환율은 1달러 당 한국 원화 단위로 표시된다.)
\begin{framedtext}
ㄱ. 명목이자율은 한국이 미국보다 높다. \par
ㄴ. 실질이자율은 한국이 미국보다 높다. \par
ㄷ. 원-달러 명목환율은 1년 간 1 \% 상승할 것으로 예상된다. \par
ㄹ. 원-달러 실질환율은 1년 간 변하지 않을 것으로 예상된다. \par
\end{framedtext}
\begin{choices}
  \item ㄱ, ㄴ
  \item ㄱ, ㄴ, ㄷ
  \item ㄱ, ㄷ, ㄹ
  \item ㄴ, ㄷ, ㄹ
  \item ㄱ, ㄴ, ㄷ, ㄹ
\end{choices}
\end{question}
\answer{3}
